\begin{center}
    {\normalsize \textbf{Giorno 3} \textbar\ Vorrei rimanere qui per sempre \textbar\ 11 Agosto 2024 \textbar\ \textbf{Tetebatu}, Lombok}
\end{center}

\noindent\rule{\textwidth}{0.4pt}
\begin{figure}[h!] % Portrait images below
    \centering
    % First portrait image (on the left)
    \begin{minipage}[h]{0.48\textwidth} % Set width equal to half the page
        \centering
        \includegraphics[width=\textwidth]{images/flags/Screenshot Tetebatu.png}
    \end{minipage}
    \hfill
    % Text
    \begin{minipage}[h]{0.48\textwidth}
        \raggedright
        \textbf{Superficie}: 4.567 km² \\
        \vspace{0.3cm}
        \textbf{Popolazione}: 3,3 milioni \\
        \vspace{0.3cm}
        \textit{L'isola}
    \end{minipage}
\end{figure}

\landscapeLargeMargin{images/day3/PXL_20240810_230650656.jpg}

\landscapeLargeMargin{images/day3/PXL_20240810_230655672.jpg}

\landscapeLargeMargin{images/day3/PXL_20240811_011914765.MP.jpg}

\landscapeLargeMargin{images/day3/PXL_20240811_015643446.MP.jpg}

\landscapeLargeMargin{images/day3/PXL_20240811_014833310.jpg}

\landscapeLargeMargin{images/day3/PXL_20240811_022739258.jpg}

\landscapeLargeMargin{images/day3/PXL_20240811_023535926.jpg}

\landscapeLargeMargin{images/day3/PXL_20240811_023609287.jpg}

\landscapeLargeMargin{images/day3/PXL_20240811_022051947.jpg}

\landscapeLargeMargin{images/day3/PXL_20240811_034604366.jpg}

\landscapeLargeMargin{images/day3/IMG_3766.JPG}

\landscapeLargeMargin{images/day3/PXL_20240811_050446006.jpg}

\landscapeLargeMargin{images/day3/PXL_20240811_091427024.MP.jpg}

\landscapeLargeMargin{images/day3/PXL_20240811_095725361.jpg}

\noindent 11 Agosto 2024 \newline

Quando mi sveglio mi sento completamente riposato. Mi sono coricato prima delle undici, sopraffatto dalla stanchezza, e ho la sensazione di aver dormito profondamente per l'intera notte, tanto che al risveglio non avverto più alcuna traccia di sonno. Guardo l'ora: è da poco passata la mezzanotte. Seguono novantadue minuti di bestemmie.

La situazione è talmente surreale da fare quasi ridere, ma stavolta non ci sono aerei da prendere l'indomani mattina e siamo reduci da una giornata di avventure, quindi sono fiducioso che riuscirò a rilassarmi e a riaddormentarmi. Infatti è proprio quello che succede, e anche se impiego un'oretta a prendere sonno alla fine ci riesco. Questa volta sono svegliato dal canto del gallo, nuovamente con la sensazione di non dover più dormire. Ricontrollo l'ora: si sono fatte le due del mattino. Vorrei buttarmi giù dalla finestra della stanza. Magari sopravviverei, ma se fossi davvero fortunato potrei atterrare sul gallo e farlo secco, e già questa sarebbe una vendetta sufficiente.

Anche Elisa viene svegliata dal canto del gallo, ma per sua fortuna si mette i tappi e riprende presto a dormire. Temo che questa volta passerò seriamente una notte in bianco, perché non c'è verso di riprendere sonno. Ripercorro tutto il mio repertorio per cercare di rilassarmi: esercizi di respirazione imparati ad apnea, pensare alla trama di un film noioso, rilassare il corpo cercando di sentire il sangue che scorre nelle vene. Niente. Verso le tre del mattino sto ripetendo a mente, in loop, l'indicativo presente del verbo essere in francese. Je suis, tu es, il est, e così via.

Dopo altre due ore insonni, proprio quando credo che la situazione non possa peggiorare, capisco finalmente perché il personale della struttura ci abbia fornito i tappi per le orecchie. Poco dopo le quattro del mattino, dalle casse della vicina moschea, giunge la voce del Muezzin che intona le preghiere a tutto volume. I galli al confronto erano una ninna-nanna. A breve le altre moschee nei dintorni seguono l'esempio, ed è così che ci troviamo nel bel mezzo di un concerto di cantilene arabe incomprensibili. Si sveglia nuovamente anche Elisa e insieme cerchiamo di infilarci i tappi così a fondo dentro le orecchie da perforarci i timpani. Si susseguono pensieri di morte e distruzione: probabilmente questa notte tradirei Ramon in cambio di un'ora di sonno.

Quando i Muezzin la smettono sono ormai le sei del mattino e si inizia a intravedere la luce dell'alba. Mi rassegno a una giornata a base di caffeina per rimanere in piedi.

Faccio però un piano per la notte successiva. Per isolarmi dai rumori esterni chiederò a Elisa di colarmi cera calda nelle orecchie da una candela accesa. Dal momento che ogni fonte di luce mi dà fastidio, mi caverò i bulbi oculari e li conserverò in un bicchiere di soluzione salina sul comodino, come se fossero la dentiera. Per ovviare alla testa che pensa sempre e non mi fa dormire dovrò trovare il modo di indurmi il coma, magari attaccandomi allo scarico di una marmitta tipo aerosol prima di spegnere la luce. La cosa più difficile di questo piano è probabilmente trovare una candela: il resto mi sembra realizzabile, sensato e soprattutto per niente eccessivo.

Mi faccio coraggio e mi alzo. Cammino dritto a fatica e mi dirigo verso la doccia. Apro l'acqua e scopro che non c'è acqua calda. In questo momento farsi una doccia fredda sarebbe come buttarsi giù dalla finestra e scoprire che invece di un gallo sono atterrato su un gatto che, come me, non riusciva a dormire e girava facendosi gli affari suoi. Chiudo l'acqua, mi vesto, preparo lo zaino per la giornata e scendo a scattare due foto al sole appena sorto.

Sono il primo turista che si aggira per la struttura a quest'ora e vengo accolto dai proprietari, in particolare dalle donne, che mi chiedono se desideri un caffè. Sì, grazie, e che sia ustionante, nero e denso come la pece. Sembrano capire anche senza che lo comunichi esplicitamente, perché la tazza che mi arriva è nera, densa, piena fino all'orlo e calda come il fuoco: dopo un inizio di giornata incerto condito dalla doccia fredda, è l'antidoto contro ogni male. Nulla potrebbe impedirmi di gustare quel caffè, a meno di trovarmi seduto al tavolo di fronte al Muezzin di stanotte: in tal caso gli colerei il caffè bollente nelle palle degli occhi e poi lo obbligherei a ingoiare la tazza tutta intera.

Sorseggiare il caffè mi fa riprendere fiducia nell'umanità. Ripensando alle preghiere musulmane fa strano vedere una popolazione praticante non araba, e mi rendo conto di come la religione qui venga vissuta davvero come culto di comunità e fratellanza, e non in maniera settoriale e fanatica come siamo abituati a vedere tramite i nostri media. È vero, abbiamo appurato che il gusto di rompere le balle al prossimo resta, ma ritengo che questa sia una caratteristica abbastanza tipica delle religioni monoteiste. Forse escludendo il buddhismo, ma con quel faccione sorridente e simpatico come si fa a essere fastidiosi?

Terminato il caffè salgo a chiamare Elisa per la colazione e come al solito ricevo una caterva di insulti perché vorrebbe dormire di più. Anch'io, principessa, anch'io, ma hai giurato "nella gioia e nel dolore" e oggi della seconda si parla, quindi trottare.

Appena scendiamo ci accolgono con un piatto di frutta esotica appena tagliata — papaia, ananas, banana e anguria — oltre al solito succo con dragon fruit e banana. Ci facciamo portare caffè, uova strapazzate per Elisa e una omelette per me. Il caffè è lo stesso di prima e con la seconda tazza iniziamo a ragionare sul serio. La colazione non è abbondante, così Eli ordina una seconda omelette, scatenando lo sconcerto delle donne in cucina: qui nessuno è abituato a consumare più di un piatto a pasto.

Terminata la colazione mi rilasso un po' sul divanetto mentre aspetto che Eli finisca di prepararsi, e nel frattempo arriva Oshi a presentarmi il programma della giornata. Prima le risaie, poi le cascate dove si può fare il bagno, infine la foresta delle scimmie. Per pranzo visiteremo un'attività locale, e al pomeriggio torneremo nel villaggio per preparare olio di cocco e macinare il caffè. Un bel piano per una lunga giornata da affrontare con due ore di sonno.

Ci incamminiamo puntuali alle nove verso le risaie, che si trovano subito al fondo del villaggio.

Oshi si rivela una guida straordinaria: conosce ogni angolo di questo posto con una precisione che lascia a bocca aperta. A occhio nudo sa indicarci da quanto tempo è stato seminato il riso, con la precisione di una settimana. Ci spiega le varie fasi di crescita del chicco — dal seme alla farina, dal latte di riso al chicco solido pronto per la raccolta — e ci fa toccare con mano ogni passaggio, fino a spremere un chicco non ancora maturo per farne uscire il latte. Quando il riso è pronto, le piante vengono sbattute con forza su un piano di legno inclinato per separarne i chicchi: un ottimo lavoro, a suo dire, per sfogarsi quando si hanno problemi nella vita, perché "you get home with the smile".

Lungo il percorso scopriamo che le risaie non sono l'unica coltura: ci sono campi di pomodori, fagiolini, soia, tabacco e soprattutto peperoncino. Qui si distinguono due varietà: il Lombok Chili, quello usato dai nativi, e il Tourist Chili o Candy Chili, più leggero e dolce, riservato ai turisti. Ci sono anche banani — che producono un solo casco per pollone e poi non ne crescono più — e alberi di cotone di vari tipi, più soffice per i vestiti e più grezzo per cuscini e materassi. Per quanto riguarda il caffè, coltivano arabica sulle montagne e robusta in pianura: la prima più acida, la seconda più amara per l'alto contenuto di caffeina.

Sul tema delle lumache, che si cibano dei raccolti e rappresentano un vero problema per gli agricoltori, la soluzione locale è tanto pratica quanto gustosa: finiscono nel saté, i tipici spiedini normalmente a base di pollo. Sul tabacco invece scopriamo che alcuni campi appartengono alla Philip Morris, mentre il resto è consumato dai locali. Il fumo è un problema serio: da una ricerca recente risulta che solo il 5% della popolazione non fuma. Non ci sono controlli nemmeno sulle strade, per cui è frequente vedere ragazzini guidare la moto senza casco e con una sigaretta in mano, magari con un passeggero. Oshi ci cita in un paio di occasioni la marijuana, ma quando gli chiedo delucidazioni mi fa il segno del tagliagole. Ricordo di aver letto che in Indonesia c'è la pena di morte per lo spaccio. Rimango nel dubbio se fosse una proposta velata o una semplice nozione, ma decido di non approfondire perché mi vedo già con una mano tagliata.

Sullo sfondo del nostro cammino si intravede il Monte Rinjani, alto 3726 metri. Oshi ci racconta il detto locale: chi scala il monte deve portare con sé solo il riso, perché il pesce lo si pesca direttamente nel lago in cima, che ne è pieno.

La conoscenza di Oshi non si limita all'agricoltura. Conosce ogni pianta medicinale del posto e ce ne mostra varie: alcune contro la diarrea e il mal di pancia — il bully belly, come lo chiama lui quando un turista si sente male — altre contro il mal di testa, altre ancora per aiutare le neo-mamme ad avere latte in abbondanza. Tra i pochi nomi che riesco a segnare c'è la Guaba, da mangiare secca contro la diarrea, e la Galanga, utile per tosse e perdita della voce, usatissima durante il COVID.

Ma Oshi ha anche un lato più giocoso. Ci prepara delle corone con delle felci che indossiamo prontamente e che ci valgono il titolo di Jungle King e Jungle Queen da parte di un'altra guida di passaggio. Ci mostra come la galanga essiccata possa essere usata come maracas, strappa degli steli e li usa come fionda per lanciare il corpo centrale come un proiettile. Infine ci realizza il jungle tattoo: appoggia una pianta sul braccio e la colpisce forte con la mano, così che rilasci una sottile polvere bianca appiccicata sulla pelle.

Una curiosità su Oshi: tra una spiegazione e l'altra, canta per tutto il tempo. La sua canzone preferita, o comunque quella che sentiamo più spesso, è Zombie dei Cranberries. È anche un grande fan della boxe e di Mike Tyson: mi chiede chi secondo me vincerebbe tra lui e Jake Paul e concordiamo entrambi che vincerebbe Tyson. Mi racconta che tramite YouTube organizzano tornei locali per promuovere la boxe e le arti marziali, con premi fino a 100 milioni di rupie, e che i migliori combattenti arrivano a competizioni nazionali e internazionali.

Il giro per le risaie dura meno di un paio d'ore ma ci riempie di nozioni dall'inizio alla fine, e sono spesso costretto a registrarlo o a prendere appunti per ricordarmi tutto. Il suo inglese non è il massimo, quindi potrei aver riportato qualche imprecisione, ma giuro: tutto in buona fede.

Terminato il tour delle risaie ci dirigiamo alle cascate sotterranee. Paghiamo l'ingresso — sessanta centesimi a testa — indossiamo le scarpette da roccia ed entriamo in una spaccatura nel suolo attraversata dall'acqua e ricoperta in alto da piante. Dopo essere passati su un fragile e umido ponte di legno, scavalchiamo a mani nude alte rocce per arrivare in un antro finale dove, con l'acqua alle ginocchia, troviamo un bacino formato da una cascata alta meno di una decina di metri ma davvero suggestiva. Oshi ci tiene gli zaini come un appendiabiti e noi andiamo a fare il bagno. L'acqua è gelida ma molto piacevole e ci immergiamo sotto il getto per qualche foto. C'è chi fa le capriole, chi si fa lo shampoo e chi posa in modo imbarazzante. Incontriamo anche il primo italiano del viaggio: un ragazzo di Biella che, dopo aver lasciato il lavoro di imbianchino, è venuto qui per due mesi e mezzo con la fidanzata fresca di dottorato.

L'acqua fredda mi ha fatto bene, perché durante il giro delle risaie ho avuto vari momenti di buio causati dalla stanchezza. Camminare sotto il sole non aiuta e a volte mi rendo conto di spegnermi completamente, rispondendo a Elisa o a Oshi con un ritardo di qualche secondo. La sensazione è quella di vedere un film con audio e video sfasati.

La fregatura vera arriva con la terza tappa: la foresta delle scimmie. Umidità più alta, percorso in salita per la maggior parte del tempo e — dettaglio che avevo rimosso — un trekking di tre quarti d'ora. In condizioni normali sarebbero stati pochi, ma oggi farei fatica a fare le scale di Via Baveno.

Prima di raggiungere la foresta facciamo tappa in un villaggio dove ci fanno assaggiare un caffè locale. Nel tragitto notiamo le loro spezie esposte per i turisti: peperoncino, lavanda, semi di caffè e molto altro. Ci sistemiamo su sedie di bambù in una stanza adibita agli ospiti e ci vengono serviti un infuso molto forte e un caffè al cacao altrettanto micidiale, oltre a chips di banana fritta che divoriamo fino all'ultima perché siamo affamati. Si tratta di coltivatori che espongono i prodotti della loro terra ai turisti come parte dei tour organizzati dalle guesthouse, così da invogliarli ad acquistare. Ci riescono anche con noi: dopo il banchetto compriamo un pacchetto di noci di macadamia e uno di fave di cacao. Fanno una più schifo dell'altra — le noci di macadamia non sanno di nulla, le fave di cacao sanno di terra. Al ritorno dovremo inventarci qualche ricetta o è facile che finiscano in pattumiera. La caffeina e la teina intanto non sono servite a nulla: mentre eravamo seduti a bere mi sarebbe bastato chiudere gli occhi per addormentarmi.

Andando verso la foresta delle scimmie passiamo davanti a dei laghetti artificiali pieni di carpe dove vengono organizzate gare di pesca. Attorno a uno di questi noto degli spot numerati per ciascun pescatore, fino a 53 attorno a un lago di una decina di metri quadrati. Chiedo come sia possibile che i pesci non finiscano subito e mi spiegano il trucco: li nutrono molto, poco prima della gara, così che non siano affamati e abbocchino a fatica. I premi sono molto alti per i locali — in alcune occasioni è stata messa in palio persino un'automobile.

Mentre ci dirigiamo verso la foresta incontriamo uomini e donne intenti a lavorare nelle risaie o a tagliare erbe, che vengono poi legate in mazzetti e portate nei villaggi per essere vendute. Incrociamo anche uno strano tipo che è salito nella foresta col motorino ed è seduto per terra con le gambe incrociate in silenzio. Proseguiamo e finalmente vediamo le prime scimmie grigie appese agli alberi. Man mano che andiamo avanti ne troviamo sempre di più e ci fermiamo a scattare foto. A un certo punto a Elisa scappa una foto col flash e tutte si dileguano spaventate.

In uno dei posti dove le guide portano i turisti scorgiamo anche due esemplari di scimmie nere, appese in cima a due altissimi alberi. Qui ci fermiamo per una pausa forzata: Oshi si adopera con un coltellino a sbucciare un bastone di canna da zucchero, regalo di un'altra guida di passaggio. Il bastone è un tubo un po' più grande di un rotolo di Scottex, con una corteccia marrone che ricorda il bambù; una volta rimossa, sotto si presenta bianco e fibroso. Se si morde questa fibra ne esce un succo zuccherino. Io provo un pezzetto e lo zucchero mi restituisce un po' di energie, ma lascio il bastone a Elisa senza assaggiare oltre perché non è particolarmente agevole da mangiare.

Riprendiamo il cammino disposti in fila indiana — Oshi davanti, io dietro, Elisa in fondo — così che io non possa vedere cosa sta facendo. Immagino stia scattando foto, ma quando mi giro vedo che ha praticamente mangiato mezza canna da zucchero, mordendo e masticando tutta la fibra bianca, e sta girando con in mano il moncherino mangiucchiato. Neanche un castoro ci sarebbe riuscito con tale velocità: a me sembrava di mangiare del legno.

Prima di lasciare la foresta troviamo un altro gruppo di scimmie grigie — stavolta almeno una dozzina — che penzolano da alberi bassi o camminano sul sentiero con i piccoli al seguito. Diventa difficile portare via Elisa, che è entrata in modalità National Geographic e scatta una foto dietro l'altra come un americano al poligono. O in un liceo. Io e Oshi ci allontaniamo e lei è costretta a seguirci per non rimanere da sola nella giungla.

Io ormai sto raggiungendo livelli di stanchezza definibili come tortura. Elisa mi fa domande e possono passare anche dieci secondi di silenzio prima che riesca ad articolare risposta. Ho sonno e fame, ma soprattutto sonno.

Almeno adesso riusciamo a placare il minore dei problemi, perché Oshi ci porta a pranzo in un'attività locale, altra meta standard dei tour organizzati e quindi piena di turisti. Il cibo però è assolutamente tradizionale: siamo affamati così prendiamo tre piatti in due e quando una bambina — probabile la figlia della proprietaria — ci si avvicina col terzo piatto è spaesata, perché anche lei non è abituata a questo mangiare abbondante dei turisti occidentali. Con un piccolo incitamento della madre alla fine ci lascia il piatto e ci augura buon appetito. In totale abbiamo ordinato pollo alla griglia con latte di cocco piccante per Elisa, riso fritto per me e noodles vegani da dividere, con pancake alla banana buonissimi per dolce, il tutto accompagnato da Coca-Cola e tè freddo. Abbiamo avuto un attimo di panico per il ghiaccio nel tè freddo, ma ci hanno assicurato che era fatto con acqua minerale, perché nemmeno loro possono bere l'acqua del rubinetto. I pancake sono stati una rivelazione: morbidi, spessi e con interi pezzi di banana gommosa all'interno.

Durante la pausa assistiamo a un paio di scene che dimostrano ancora di più l'umiltà degli indonesiani. Quando un turista tedesco offre al cameriere una mancia, quest'ultimo rimane interdetto e non capisce cosa voglia quello straniero con una mazzetta di banconote in mano: occorre dirgli "for you" per convincerlo ad accettare. In un altro caso, un turista che lascia la mancia dopo aver pagato alla cassa viene inseguito fino fuori dalla cassiera che vuole restituirgli i soldi, e deve insistere perché li accetti.

Finito il pranzo segue un attimo di relax su una delle loro piattaforme di bambù. Chiudo gli occhi e in un minuto sono già in dormiveglia, poi mi sveglio e scambio quattro chiacchiere con una guida locale con un occhio solo. A un certo punto Oshi ci dice che è giunto il momento di tornare, arriva un pickup e ci carica sul cassone: ci mettiamo in piedi dietro l'abitacolo e percorriamo il tragitto come il miglior team di giardinieri messicani.

Giunti ad Al Sasaki chiedo un'ora di pietà a Elisa e Oshi, perché altrimenti non arriverò in piedi fino a sera. Mentre salgo le scale pregusto quell'ora di sonno come un carcerato la libertà. Mi siedo sul letto, mi tolgo le scarpe e proprio quando sto per coricarmi il Muezzin della vicina moschea parte a tutto volume con la serenata. Non ci voglio credere. Elisa si mette a ridere dell'assurdità della scena. Io sono disperato, ma talmente assente che mi corico, metto tappi e mascherina e mi addormento in pochi secondi nonostante il casino. Forse, ora che ci ripenso, un bestemmino l'ho tirato, ma non ci giurerei.

Quando Elisa mi sveglia dopo un'ora ricordo a malapena di essere in Indonesia.

A fatica mi alzo e scendiamo nell'area comune, dove ci aspettano per preparare l'olio di cocco. Ci dirigiamo al fondo del villaggio con le donne di casa e Oshi, intento a sbucciare una noce di cocco con un machete. Con abili colpi riesce a togliere la corteccia esterna in pochi minuti, con la precisione di chi ha già ripetuto lo stesso gesto centinaia, se non migliaia di volte. Ci consegna poi il cocco fresco così che io ed Elisa possiamo grattugiarlo in un recipiente. Dopo vari sforzi abbiamo sbriciolato tutto il cocco e io sono riuscito a tagliarmi un dito. Il cocco viene mischiato con acqua e poi strizzato; la polpa che ne esce viene messa in un pentolino a bollire sul fuoco, alimentato con la buccia più esterna del cocco. Ogni gesto ci viene mostrato con precisione e siamo poi noi a metterlo in pratica. Cuocendo, la polpa separa l'acqua dal concentrato di cocco che rimane in superficie; questo viene raccolto con un mestolo e messo a friggere su una padella. Dalla frittura si separa l'olio di cocco, che viene filtrato, mentre la parte solida diventa croccante e ci viene data da mangiare come snack. Non si butta via nulla. L'olio viene infine colato in una bottiglietta di vetro che ci consegnano come souvenir.

Inizia poi la seconda parte dell'attività: la tostatura del caffè. Prima si tosta il riso, poi vi si aggiungono i chicchi di caffè e il tutto viene ritostato insieme. Il riso serve ad allungare il caffè e renderlo meno forte, necessario perché qui ne bevono tantissimo durante il giorno. Quando sono pronti, riso e caffè vengono posti in un contenitore di legno e macinati pestando con dei grossi bastoni. Io ed Eli ci cimentiamo nella macinazione, ma non siamo bravi quanto le donne di casa. A fine lavorazione ci consegnano una bustina con il nostro caffè, non prima di averne usata una parte per prepararcene una tazza calda: il fatto di averlo prodotto a mano dalla materia prima lo rende uno dei caffè migliori che abbia mai assaggiato. O più probabilmente sono ancora molto stanco e l'ennesima tazza della giornata è un toccasana.

Nello svolgere queste attività siamo sempre stati seguiti dallo sguardo attento ma discreto di Ronnie, il proprietario, che controlla la situazione rimanendo in disparte. È insieme a un amico che chiacchiera e banchetta, e quando può ci offre da mangiare. Dapprima rifiutiamo, ma vista l'insistenza ci lasciamo convincere. La prima cosa che assaggiamo è un dolce realizzato con cocco, farina e zucchero di canna: ha una consistenza gelatinosa molto densa e un sapore buonissimo, tanto che ne prendo quattro pezzi compresi quelli di Elisa, che non apprezza quanto me. La seconda offerta è frutta fresca tagliata a pezzi da immergere in una salsa piccante fatta da loro. Gli chiedo se si tratta del Lombok chili o del Tourist chili, per capire il grado di piccante e per fare vedere che abbiamo imparato qualcosa dalla giornata con Oshi. Nonostante si tratti del chili piccante assaggiamo lo stesso e anche qui rimaniamo stupiti dal sapore. Io ne prendo a volontà.

Tutti questi spuntini ci aprono lo stomaco e dopo il caffè arriva il momento di cenare. Siamo molto stanchi e sarebbe bello goderci semplicemente la cena cucinata da loro, ma quando Oshi ci chiede "Wanna cook or wanna relax?" rispondiamo senza pensarci: "Cook". Come al solito ci piace farci del male.

Elisa ha pensato tutto il giorno ai noodles alle verdure e li riordina, mentre io voglio provare un piatto nuovo tipico di Lombok e scelgo il pollo con riso, verdure e latte di cocco. Cucino insieme alle donne di casa e mentre Elisa accompagna Oshi a cuocere il pollo, assisto a una scena che dimostra la massima cura di queste persone verso i loro ospiti: una delle donne più giovani raccoglie con un mestolo un pezzetto di verdura per tipo dalla pentola e le consegna alla donna più anziana, che le assaggia una per una con attenzione e meticolosità per verificare se siano cotte. Solo allora dà il via libera per aggiungere la salsa di soia agrodolce. A questo punto mi lancio nel mio debole show da spadellatore e tutte mi chiamano chef. In qualsiasi altro contesto l'avrei considerata una presa in giro, ma qui è evidente che sono sincere. Quasi mi commuovono.

Elisa intanto è con Oshi che griglia il pollo dopo averlo caramellato con latte di cocco, il tutto su una brace fatta di gusci di cocco. Mi fa notare che tutti ci guardano straniti: io sono rimasto in cucina ai fornelli con le donne del villaggio mentre lei è andata con gli uomini a cuocere la carne. Qui le donne non fanno i lavori degli uomini e viceversa, fatta eccezione per Oshi che è un tuttofare esperto.

Mentre prepariamo il pollo, Ronnie sta cenando con un suo compare sulla piattaforma di bambù e in continuazione ci chiede se abbiamo bisogno di qualcosa, offrendoci le sue chips di riso fatte in casa. Anche se diciamo di no, alla fine se le toglie per darle a noi, convinto che siano buone col pollo. Ogni scusa è buona per essere gentili.

Quando tutto è pronto ci sediamo a gustare la cena ed è tutto delizioso come la sera precedente. Ci spiace davvero lasciare questo angolo di paradiso, ma abbiamo un'altra avventura che ci aspetta l'indomani.

\begin{center}
    % Decorative line
    \noindent\rule{0.6\textwidth}{0.4pt}
\end{center}

Apro una piccola parentesi di sfogo: Elisa qui non ci voleva venire. Quando le ho detto che saremmo andati dai nativi si immaginava di dormire per terra in una capanna, di mangiare con le mani da piatti sporchi, di non capire niente perché nessuno parla italiano, e gne gne gne. A nulla sono valse le mie spiegazioni e rassicurazioni: aveva deciso che avremmo preso qualche malattia e saremmo morti di fame. Quando a fine serata pronuncia la frase "Vorrei rimanere qui per sempre", per poco non le lancio una gallina. Si è innamorata del posto, l'inglese non è stato un grosso problema — anzi l'ha spronata a cimentarsi in autonomia — ha mangiato bene come non avrebbe mai immaginato e si è ricreduta su ogni pregiudizio che aveva. Non sui nativi di Lombok, ma sul mio planning. Lo scrivo perché se stamperò questo diario potrò autocitarmi quando mi servirà dire "Te l'avevo detto". Fine parentesi.

\begin{center}
    % Decorative line
    \noindent\rule{0.6\textwidth}{0.4pt}
\end{center}

Diamo la buonanotte a tutti e saliamo in stanza. Prima di andare a dormire decido di corazzarmi: prendo la bellezza di quarantasette gocce di Quietil, medicinale omeopatico per dormire. La dose consigliata sarebbe stata trenta, io puntavo a sessanta, ma Elisa mi controllava perché temeva andassi in coma e mi sono dovuto fermare prima. Non c'è fretta di fare i bagagli la sera così andiamo subito a letto, consapevoli che verremo svegliati dal Muezzin in piena notte e sperando di accumulare un po' di sonno prima.


\pagestyle{empty} % disable numbers

\twoImagesLargeMargins{images/day3/PXL_20240810_230815381.jpg}
                      {images/day3/PXL_20240810_230929822.jpg}

% coppia
\landscapeLargeMargin{images/day3/PXL_20240810_231614099.jpg}
\landscapeLargeMargin{images/day3/PXL_20240811_000725026.jpg}
%

%coppia
\landscapeLargeMargin{images/day3/PXL_20240811_012257649.MP.jpg}
\twoImagesLargeMargins{images/day3/PXL_20240811_012730414.MP.jpg}
                      {images/day3/PXL_20240811_013512953.jpg}
%

%coppia
\portraitFullscreen{images/day3/IMG_3198.JPG}
\landscapeThinMarginCentered{images/day3/IMG_3260.JPG}
%

%coppia
\landscapeLargeMargin{images/day3/PXL_20240811_020412829.jpg}
\landscapeLargeMargin{images/day3/PXL_20240811_020441511.jpg}
%

%coppia
\landscapeLargeMargin{images/day3/PXL_20240811_015755369.jpg}
\landscapeLargeMargin{images/day3/PXL_20240811_020106485.MP.jpg}
%

%coppia
\landscapeLargeMargin{images/day3/PXL_20240811_020841057.jpg}
\landscapeLargeMargin{images/day3/PXL_20240811_023638421.MP.jpg}
%

%coppia
\splitImageLeft{images/day3/IMG_3317.JPG}
\splitImageRight{images/day3/IMG_3317.JPG}
% 

%coppia
\twoImagesThinMargins{images/day3/IMG_3730.JPG}
                      {images/day3/IMG_3371.JPG}
\twoImagesThinMargins{images/day3/IMG_3531.JPG}
                      {images/day3/IMG_3480.JPG}
%

%coppia
\twoImagesLargeMargins{images/day3/IMG_3443.JPG}
                      {images/day3/PXL_20240811_035414133.jpg}
\twoImagesLargeMargins{images/day3/IMG_3426.JPG}
                      {images/day3/IMG_3430.JPG}
%

%coppia
\landscapeThinMargin{images/day3/IMG_3453.JPG}
\twoImagesLargeMargins{images/day3/IMG_3444.JPG}
                      {images/day3/IMG_3450.JPG}
%

%coppia
\fourPortraitImages{images/day3/PXL_20240811_025138114.jpg}
                   {images/day3/PXL_20240811_025419563.jpg}
                   {images/day3/PXL_20240811_025445044.MP.jpg}
                   {images/day3/PXL_20240811_025922590.jpg}
\twoImagesLargeMargins{images/day3/PXL_20240811_025944014.jpg}
                      {images/day3/PXL_20240811_025509136.jpg}
%
 
%coppia
\splitImageLeft{images/day3/IMG_3394.JPG}
\splitImageRight{images/day3/IMG_3394.JPG}
%

%coppia
\landscapeThinMargin{images/day3/PXL_20240811_034603976.jpg}
\twoImagesLargeMargins{images/day3/PXL_20240811_034652563.jpg}
                      {images/day3/PXL_20240811_034746227.jpg}
%

%coppia
\landscapeThinMarginCentered{images/day3/IMG_3510.JPG}
\twoImagesLargeMargins{images/day3/IMG_3505.JPG}
                      {images/day3/IMG_3527.JPG}
%

%coppia
\landscapeThinMarginCentered{images/day3/PXL_20240811_052157140.MP.jpg}
\landscapeThinMarginCentered{images/day3/PXL_20240811_055627220.jpg}
%

%coppia
\twoImagesLargeMargins{images/day3/IMG_3624.JPG}
                      {images/day3/IMG_3631.JPG}
\threeImagesLargeMarginsInverted{images/day3/IMG_3662.JPG}
                                {images/day3/IMG_3652.JPG}
                                {images/day3/IMG_3655.JPG}
                                
%

%coppia
\fourPortraitImages{images/day3/PXL_20240811_090832711.NIGHT.jpg}
                   {images/day3/PXL_20240811_090919889.NIGHT.jpg}
                   {images/day3/PXL_20240811_091013098.NIGHT.jpg}
                   {images/day3/PXL_20240811_091225227.MP.jpg}
\landscapeLargeMargin{images/day3/PXL_20240811_091638592.MP.jpg}
%

%coppia
\landscapeThinMargin{images/day3/PXL_20240811_091730612.MP.jpg}
\twoImagesLargeMargins{images/day3/PXL_20240811_091427642.MP.jpg}
                      {images/day3/PXL_20240811_092659694.jpg}
%

%coppia
\twoImagesThinMargins{images/day3/PXL_20240811_095255229.MP.jpg}
                     {images/day3/PXL_20240811_094559918.MP.jpg}
\eightLandscapeImages{images/day3/PXL_20240811_093521357.jpg}
                     {images/day3/PXL_20240811_093547434.MP.jpg}
                     {images/day3/PXL_20240811_093652486.MP.jpg}
                     {images/day3/PXL_20240811_093856691.jpg}
                     {images/day3/PXL_20240811_094649477.jpg}
                     {images/day3/PXL_20240811_094601918.MP.jpg}
                     {images/day3/PXL_20240811_093939515.MP.jpg}
                     {images/day3/PXL_20240811_095106418.jpg}
%
\twoImagesLargeMargins{images/day3/PXL_20240811_095321675.jpg}
                      {images/day3/PXL_20240811_095517063.MP.jpg}
\twoImagesLargeMargins{images/day3/PXL_20240811_095725361.jpg}                      
                      {images/day3/PXL_20240811_095333494.MP.jpg}

%coppia
\landscapeLargeMargin{images/day3/PXL_20240811_100027500.MP.jpg}
\landscapeLargeMargin{images/day3/PXL_20240811_100115657.MP.jpg}
%

%coppia
\landscapeLargeMargin{images/day3/PXL_20240811_100907195.jpg}
\twoImagesThinMargins{images/day3/PXL_20240811_100240617.jpg}
                     {images/day3/PXL_20240811_101059807.jpg}
%

%coppia
\twoImagesThinMargins{images/day3/PXL_20240811_101403930.jpg}
                     {images/day3/PXL_20240811_101503552.jpg}
\landscapeThinMarginCentered{images/day3/PXL_20240811_122342044.MP.jpg}
%

%coppia
\landscapeLargeMargin{images/day3/PXL_20240811_130259070.jpg}

% rimuovere la prossima se non serve a pareggiare le pagine
% \landscapeLargeMargin{images/day3/PXL_20240811_130328625.MP.jpg}
%

\pagestyle{fancy} % enable numbers